% -----------------------------------------------------------
\section{Salvation}
% -----------------------------------------------------------
	\label{SALVATION}
	
% % ----------------------
% \subsection{See also}
% % ----------------------

% % ----------------------
% \subsection{1828 American Dictionary of the English Language} % *1828
% % ----------------------
	% \label{SALVATION-1828}

	% \begin{compactitem}
		% \item
	% \end{compactitem}

% % ----------------------
% \subsection{Oxford English Dictionary} % *OED
% % ----------------------
	% \label{SALVATION-OED}

	% \begin{compactitem}
		% \item[]
	% \end{compactitem}
	
% % ----------------------
% \subsection{Buck's Theological Dictionary (1833)} % *BTD
% % ----------------------
	% \label{SALVATION-BTD}

	% \begin{compactdesc}
		% \item[PAGE\#]
	% \end{compactdesc}
	
% % ----------------------
% \subsection{Restoration Lexicon} % *RL *RESTORATION LEXICON
% % ----------------------
	% \label{SALVATION-OED}

	% \begin{compactitem}
		% \item[]
	% \end{compactitem}

% ----------------------
\subsection{Scriptural Usage} % *SCRIPTURES
% ----------------------
	\label{SALVATION-SCRIPTURES}
	
	% \begin{tabular}{ll}
		% Total occurrences of ``'' in the scriptures & 0
		% Total occurrences of ``'' in the scriptures & 0
	% \end{tabular}

	% % -----
	% \subsubsection{Old Testament} % *OT
	% % -----
		% \label{SALVATION-OT}
	
		% % --
		% \paragraph{KJV}
		% % --
			% \label{SALVATION-OT-KJV}
	
			% \begin{tabular}{ll}
				% Total occurrences in the OT & 
			% \end{tabular}
			
			% \begin{compactdesc}
				% \item[]
			% \end{compactdesc}
			
		% % --
		% \paragraph{Hebrew Bible}
		% % --
			% \label{SALVATION-HEBREW}
		
			% %
			% \subparagraph{\texorpdfstring{\he{sample word}}{}} % paste actual hebrew text here
			% %
				% \label{PAGA} % whatever is in step bible without periods
			
				% \begin{compactdesc}
					% \item[HALOT , PDF ] \textbf{} % Use the 1994 PDF
						% \begin{compactitem}
							% \item[1]
						% \end{compactitem}
					% \item[BDB , PDF ] \textbf{}
						% \begin{compactitem}
							% \item[1]
						% \end{compactitem}
					% \item[DCH PDF ] \textbf{} % don't include the actual page number, they repeat from volume to volume
						% \begin{compactitem}
							% \item[1]
						% \end{compactitem}
				% \end{compactdesc}
				
				% \begin{compactdesc}
					% \item[Some OT uses] \textbf{}
						% \begin{compactdesc}
							% \item[]
						% \end{compactdesc}
				% \end{compactdesc}
			
		% % --
		% \paragraph{Septuagint}
		% % --
			% \label{SALVATION-LXX}
		
			% %
			% \subparagraph{\texorpdfstring{\gr{sample word}}{}}
			% %
		
	% % -----
	% \subsubsection{New Testament} % *NT
	% % -----
		% \label{SALVATION-NT}
	
		% % --
		% \paragraph{KJV}
		% % --
			% \label{SALVATION-NT-KJV}		
	
			% \begin{tabular}{ll}
				% Total occurrences in the NT & 
			% \end{tabular}
			
			% \begin{compactdesc}
				% \item[]
			% \end{compactdesc}
			
		% % --
		% \paragraph{Greek New Testament} % *GREEK
		% % --
			% \label{SALVATION-GREEK}
		
			% %
			% \subparagraph{\texorpdfstring{\gr{sample word}}{}}
			% %
				% \label{KATALLAGE}
			
				% \begin{compactdesc}
					% \item[BDAG ] \textbf{}
						% \begin{compactitem}
							% \item 
						% \end{compactitem}
					% \item[CGL , PDF ] \textbf{}
						% \begin{compactitem}
							% \item 
						% \end{compactitem}
					% \item[LSJ , PDF ] \textbf{}
						% \begin{compactitem}
							% \item 
						% \end{compactitem}
					% \item[Brill , PDF ] \textbf{}
						% \begin{compactitem}
							% \item 
						% \end{compactitem}
				% \end{compactdesc}
				
				% \begin{tabular}{ll}
					% Total occurrences in the NT & 
				% \end{tabular}
				
				% \begin{compactdesc}
					% \item[Some NT uses] \textbf{}
						% \begin{compactdesc}
							% \item[]
						% \end{compactdesc}
				% \end{compactdesc}
				
			% %
			% \subparagraph{TDNT: \texorpdfstring{\gr{sample word}}{}  (3:536--556)}
			% %
				% \label{KATALLAGE-TDNT}
		
	% % -----
	% \subsubsection{Book of Mormon} % *BOM
	% % -----
		% \label{SALVATION-BOM}
	
		% \begin{tabular}{ll}
			% Total occurrences in the BOM & 
		% \end{tabular}
		
		% \begin{compactdesc}
			% \item[]
		% \end{compactdesc}
		
	% % -----
	% \subsubsection{Doctrine and Covenants} % *DC
	% % -----
		% \label{SALVATION-DC}
	
		% \begin{tabular}{ll}
			% Total occurrences in the DC & 
		% \end{tabular}
		
		% \begin{compactdesc}
			% \item[]
		% \end{compactdesc}
		
	% % -----
	% \subsubsection{Pearl of Great Price} % *PGP
	% % -----
		% \label{SALVATION-PGP}
	
		% \begin{tabular}{ll}
			% Total occurrences in the PGP & 
		% \end{tabular}
		
		% \begin{compactdesc}
			% \item[]
		% \end{compactdesc}
		
% ----------------------
\subsection{Key Phrases} % *KEYPHRASES *PHRASES
% ----------------------
	\label{SALVATION-PHRASES}

	\begin{compactdesc}
	
		\item[temporal salvation] \textbf{2} | \hyperref[DC89-2]{DC 89:2}; \hyperref[MOSES-7-42]{Moses 7:42}
			\begin{compactdesc}
				\item[Bible anchor verses] ---
					% \begin{compactdesc}
						% \item[Romans 5:5] \gr{}
					% \end{compactdesc}
				% \item[Commentary] ---
			\end{compactdesc}
			
	\end{compactdesc}
	
% % ----------------------
% \subsection{Bible Dictionaries} % *DICTIONARIES
% % ----------------------
	% \label{SALVATION-BD}

% ----------------------
\subsection{Restoration Usage} % *RESTORATION
% ----------------------
	\label{SALVATION-RESTORATION}

	% -----
	\subsubsection{Joseph Smith} % *JS
	% -----
		\label{SALVATION-JS}
	
		\begin{compactdesc}
		
			\item[{\hyperlink{EMS-1843-JS-14-MAY-1843-A-WW-a}{14 May 1843 WW-a}}] The principl of knowledge is the principle of Salvation the Principle can be comprehended, for any one that cannot get knowledge to be Saved will be damned. The Principl of Salvation is given to us through the knowledge of Jesus Christ Salvation is nothing more or less than to triumph over all our enemies \& put them under our feet \& when we have power to put all enemies under our feet in this world \& a knowledge to triumph over all evil spirits in the world to come then we are saved as in the case of Jesus he was to reign untill he had put all enemies under his feet \& the last enemy was death
				\begin{compactitem}
					\item For this idea of ``salvation'' as triumphing over enemies, see \hyperref[LUKE-1-69]{Luke 1:69--71}, 
				\end{compactitem}
			
			\item[{\hyperref[EMS-1843-JS-21-MAY-1843-WR]{21 May 1843 (Willard Richards)}}] \uline{salvation}. is for a man to be seveed from all his enemies.--- until a man can triumph over death. he is not saved. knowlidge will do this.
				\begin{compactitem}
					\item At ``seveed,'' a JSP footnote says: ``\textsc{text}: Likely `seve[r]ed' but possibly `saveed [saved]'.''
				\end{compactitem} 
			
			\item[{\hyperref[EMS-1843-JS-21-MAY-1843-HC]{21 May 1843 (Howard Coray)}}] What is Salvation. Salvation is for a man to be Saved from all his enemies even our last enemy which is death
			
		\end{compactdesc}
		
	% % -----
	% \subsubsection{Temple Ordinances}
	% % -----
		% \label{SALVATION-TEMPLE}
	
		% \begin{compactdesc}
			% \item[{\hyperref[KEY]{Date/Source}}]
		% \end{compactdesc}
	
	% % -----
	% \subsubsection{Brigham Young} % *BY
	% % -----
		% \label{SALVATION-BY}
	
		% \begin{compactdesc}
			% \item[{\hyperref[KEY]{Date/Source}}]
		% \end{compactdesc}
	
% ----------------------
\subsection{Commentary and Studies} % *COMMENTARY *STUDIES
% ----------------------
	\label{SALVATION-COMMENTARY}

	% % -----
	% \subsubsection{Key Points}
	% % -----
		% \label{SALVATION-KEYS}
	
		% \begin{compactitem}
			% \item
		% \end{compactitem}
		
	% -----
	\subsubsection{Universal salvation in the restored gospel}
	% -----
		\label{SALVATION-universal}
		
		\begin{compactitem}
			\item \hyperref[DC52-15]{DC 52:15--18}
			\item \hyperref[DC58-7]{DC 58:7--11}
			\item \hyperref[DC76-44]{DC 76:44}, thinking about the sons of perdition, but it also speaks of the ``end'' of their punishment; think of \hyperref[DC19]{DC 19} stuff there, and also \hyperref[NATIVEELEMENT]{Native Element} stuff too
		\end{compactitem}